\SourceLeaf{244}%
1830 zu finden sein möchte. Dagegen aber kann ich Ihnen den Titel einer andern Schrift nachweisen, die Sie ohne Zweifel sehr leicht durch den Buchhandel erhalten können, und die nur 4 Bogen stark ist:

\begin{center}
»Geometrische Untersuchungen zur Theorie der Parallel-Linien\\
von \textsc{Nicolaus} \textsc{Lobatschewsky} Kais. Russischem Staatsrath \emph{etc.}\\
Berlin 1840 in der G. Finckeschen Buchhandlung.«
\end{center}

Ich erinnere mich, in \textsc{Gersdorfs} Repertorium damals eine sehr wegwerfende Recension dieses Buchs gelesen zu haben, die (nemlich die Recension) übrigens für jeden etwas kundigen Leser das Gepräge hatte, von einem ganz unwissenden Menschen herzurühren. Seitdem ich Gelegenheit gehabt habe, diese kleine Schrift selbst einzusehen, muss ich ein \textls{sehr vortheilhaftes} Urtheil darüber fällen. Namentlich hat sie viel mehr Concinnität und Präcision, als die grössern Aufsätze des \textsc{Lobatschewsky}, die mehr einem verworrenen Walde gleichen, durch den es, ohne alle Bäume erst einzeln kennen gelernt zu haben, schwer ist, einen Durchgang und Übersicht zu finden.

Über die \textsc{Crelle} 17 p. 303 angeführte experimentelle Begrenzung habe ich aber nichts in der Schrift von 1840 gefunden und ich werde mich daher wohl entschliessen müssen, einmal deswegen an Hrn. \textsc{Lobatschewsky} selbst zu schreiben, dessen Aufnahme als Correspondent unserer Societät ich vor einem Jahre veranlasst habe. Vielleicht schickt er mir dann den Kasanschen Boten. Doch wäre es möglich, dass sich in den folgenden Jahrgängen der Kasanschen gelehrten Schriften von 1836---1838, wo auch lange Aufsätze von \textsc{Lobatschewsky} stehen, etwas darüber findet. Ich besitze diese zwar, habe aber bisher, aus dem in meinem vorigen Briefe erwähnten Grunde, mich bisher nicht näher mit ihnen bekannt gemacht.

In seinem Danksagungsschreiben wegen Aufnahme in die Societät schrieb mir übrigens \textsc{Lobatschewsky} damals, dass seine vielen administrativen Geschäfte (er scheint Rector perpetuus der Universität zu sein) ihn jetzt aus den wissenschaftlichen Arbeiten ganz herausgebracht hätten.

Für heute schliesse ich mit der Bitte bald wieder mit einigen Zeilen zu erfreuen

\begin{flushright}
Ihren treu ergebenen\\
\textsc{C. F. Gauss}.
\end{flushright}

\par\bigskip
